\documentclass{article}
\usepackage[utf8]{inputenc}
\usepackage[T1]{fontenc}
\usepackage{geometry}
\usepackage{graphicx}
\usepackage{float}
\usepackage{hyperref}
\usepackage[english,indonesian]{babel}

\title{Optimasi Pemrosesan Template LaTeX Menggunakan Pendekatan Berbasis AI}
\author{Academic Support Agent}
\date{\today}

\begin{document}
\maketitle

\begin{abstract}

\end{abstract}

\chapter{PENDAHULUAN}


\chapter{TINJAUAN PUSTAKA}


% Original content (commented out):
% \documentclass[a4paper, 11pt, titlepage]{article}

% Set page size and margins
% Replace `letterpaper' with`a4paper' for UK/EU standard size
\usepackage[a4paper,top=2.25cm,bottom=2.25cm,left=2cm,right=2cm,marginparwidth=2cm]{geometry}

% Useful packages
\usepackage{amsmath}
\usepackage{graphicx}
\usepackage[toc,page]{appendix}
\usepackage[colorlinks=true, allcolors=blue]{hyperref}
\usepackage{caption}
\usepackage{titlepic}
\usepackage{titling}

% Multi column
% \usepackage{multicol}
\usepackage{wrapfig}
\usepackage{stfloats}
\usepackage{placeins}

% Bibliography
\usepackage[urldate=iso,seconds=true]{biblatex} %Imports biblatex package
\addbibresource{sample.bib} %Import the bibliography file

%------------------- VARIABLES -------------------%
\def\majorheading{MAJOR HEADING} % e.g. Unit name: EGH400-1 - Research Project 1
\def\minorheading{MINOR HEADING} % e.g. Group or type of report: Group 42, or Progress Report
\def\reporttitle{REPORT TITLE} % e.g. My Project Title, or ROS2 on a PenguinPi
\def\reportsubtitle{SUB TITLE} % (OPTIONAL - comment out to remove) e.g. Modernising Educational Robotics
\def\institutionname{Queensland University of Technology}

% Authors. If there are multiple, separate by \and commands. Use \\ to separate lines within a single author's entry -- for example, to give the author's institution or address. 
\def\authorname{AUTHOR 1 \and AUTHOR 2} % Can include student number, e.g. First Last - N123456789

\title{\reporttitle}
\author{\authorname}
\date{\today}

%-------------------------------------------------%

% Headers and footers
\usepackage{fancyhdr}
\pagestyle{fancy}
\fancyhf{}
\lhead{\reporttitle}
\rhead{\majorheading \space-\space \minorheading}
\setlength{\headheight}{13.6pt}
\addtolength{\topmargin}{-1.6pt}
\rfoot{Page \thepage}

\begin{document}
\begin{titlepage} % Suppresses displaying the page number on the title page and the subsequent page counts as page 1
\newcommand{\HRule}{\rule{\linewidth}{0.5mm}} % Defines a new command for horizontal lines, change thickness here

\center % Centre everything on the page

%------------------------------------------------
%	Headings
%------------------------------------------------

\textsc{\LARGE \institutionname}\\[1.5cm] % Main heading such as the name of your university/college

\textsc{\Large \majorheading}\\[0.5cm] % Major heading such as course name

\textsc{\large \minorheading}\\[0.5cm] % Minor heading such as course title

%------------------------------------------------
%	Title
%------------------------------------------------

\HRule\\[0.8cm]

{\huge\bfseries \reporttitle}\\[0.4cm] % Title of your document
\ifx\reportsubtitle\undefined 
\else
    {\textsc{\Large \reportsubtitle}}\\[0.4cm] % Subtitle
\fi

\HRule\\[1.5cm]

%------------------------------------------------
%	Author(s)
%------------------------------------------------

% {\large\textit{Authors}}\\
\theauthor \\[0.5cm]

%------------------------------------------------
%	Date
%------------------------------------------------

\vfill\vfill\vfill % Position the date 3/4 down the remaining page

{\large\today} % Date, change the \today to a set date if you want to be precise

%------------------------------------------------
%	Logo
%------------------------------------------------

\vfill\vfill
\includegraphics[width=0.2\textwidth]{images/qut-logo.jpg}\\[1cm] % Include a department/university logo - this will require the graphicx package
 
%----------------------------------------------------------------------------------------

\vfill % Push the date up 1/4 of the remaining page
	
\end{titlepage}
% \maketitle
{\hypersetup{linkcolor=black} 
% \let\clearpage\relax 
\thispagestyle{empty}
\tableofcontents
\listoffigures
\listoftables
\clearpage
}
\pagenumbering{arabic}

%------------------------------------------------
%	Reference Latex Functions
%------------------------------------------------

% ----------- Section header -----------
% \section{Importing the data}

% ----------- Section body -----------
% The first step taken in the data analysis was to import the dataset then perform preliminary checks on the data to make sure that it was imported correctly, checking that the data was of the correct type (string, numeric, date etc). As the dataset was provided in a CSV file format, the data was imported into MATLAB with the \verb_readtable_ command. Once the dataset had been imported, the first and last 100 rows of data where viewed to check the data types of the columns (Table \ref{tab:data_type}). The dataset was also observed to have 1,000,000 rows.

% ----------- Reference figures, tables, etc. with \ref (see above) -----------

% ----------- Cite bibliography source with \cite -----------
% \cite{greenwade93}

% ----------- Table -----------
% \begin{table}[h]
%     \centering
%% ----------- Table columns and sizes -----------
%     \begin{tabular}{|p{0.2\linewidth}|p{0.4\linewidth}|p{0.24\linewidth}|} 
%     \hline
%% ----------- Column names -----------
%     {Column} & {Assigned Name} & {Data type} \\ 
%     \hline
%% ----------- Column data seperated by &, line ended in \\ -----------
%     One & IV 1 & Datetime \\
%     Two & IV 2 & Numeric \\
%     Three & IV 3 & Numeric \\
%     Four & IV 4 & Numeric \\
%     Five & DV 1 or Categories & String \\
%     Six & DV 2 & Datetime \\
%     \hline
%     \end{tabular}
%     \caption{Column data types and assigned names}
%% ----------- Table label (used for references) -----------
%     \label{tab:data_type}
% \end{table}

% ----------- Figure -----------
%% For figure placement, use the following codes on \begin{figure*}[b]. b - bottom of page, t - top of page, h - here (inline), can combine in order of preference, e.g. ht prefer here then top. Can add ! to "insist" on placement, e.g. !htb.
% \begin{figure*}[b] 
% \centering
% \includegraphics[width=0.9\linewidth]{images/qut-logo.jpg}
% \caption{Caption}
% \label{fig:my_figure}
% \end{figure*}

% ----------- Verbatim environment used for code text ----------- 
% \begin{verbatim}
% \end{verbatim}

% ----------- Inline Maths ----------- 
% \LaTeX{} is great at typesetting mathematics. Let $X_1, X_2, \ldots, X_n$ be a sequence of independent and identically distributed random variables with $\text{E}[X_i] = \mu$ and $\text{Var}[X_i] = \sigma^2 < \infty$, and let
% $$S_n = \frac{X_1 + X_2 + \cdots + X_n}{n}
%       = \frac{1}{n}\sum_{i}^{n} X_i$$
% denote their mean. Then as $n$ approaches infinity, the random variables $\sqrt{n}(S_n - \mu)$ converge in distribution to a normal $\mathcal{N}(0, \sigma^2)$.

% ----------- Multiline equations (aligned) ----------- 
%% Align single equation
% \begin{equation} \label{eq1}
% \begin{split}
% A & = \frac{\pi r^2}{2} \\
%  & = \frac{1}{2} \pi r^2
% \end{split}
% \end{equation}

%% Aligning several equations
% \begin{align*} 
% 2x - 5y &=  8 \\ 
% 3x + 9y &=  -12
% \end{align*}

%% Aligning several equations (multiple columns)
% \begin{align*}
% x&=y           &  w &=z              &  a&=b+c\\
% 2x&=-y         &  3w&=\frac{1}{2}z   &  a&=b\\
% -4 + 5x&=2+y   &  w+2&=-1+w          &  ab&=cb
% \end{align*}

%% Grouping and centering equations
% \begin{gather*} 
% 2x - 5y =  8 \\ 
% 3x^2 + 9y =  3a + c
% \end{gather*}

% ----------- Lists ----------- 
% You can make lists with automatic numbering \dots. Start typing \enumerate or \itemize, then tab to autocomplete for a shortcut.
% \begin{enumerate}
% \item Like this,
% \item and like this.
% \end{enumerate}
% \dots or bullet points \dots
% \begin{itemize}
% \item Like this,
% \item and like this.
% \end{itemize}

%------------------------------------------------
%	Main Report
%------------------------------------------------

\section{Introduction}

\section{Body}
\subsection{Subsection 1}
\subsection{Subsection 2}

\section{Conclusion}

%% -----------  Insert references (bibliography) ----------- 
\printbibliography

%% -----------  Appendices ----------- 
% \onecolumn
% \begin{appendices}

% \end{appendices}

\end{document}

\end{document}